Konvolūcijas slāņi ģenerē pazīmju kartes, iegūstot vizuālas pazīmes, kas līdzinās orientētu malu detektoriem un Gabora filtru
izejām \cite{tavanaei2019}. Konvolūcijas slāņus plaši izmanto \Gls{snn} arhitektūrās, parasti aizstājot \Gls{relu}
aktivācijas ar \Gls{if} vai \Gls{lif} neironiem \cite{cao2015}. Šādi tīkli objektu detektēšanā notikumu attēlu datu kopās
pārspēj tradicionālos \Gls{ann} \cite{zhang2023}. Ir izveidoti \Gls{snn}, kuru pamatā ir esošas \Gls{cnn} arhitektūras,
piemēram, \textit{VGG} \cite{na2022}. Dziļākie slāņi efektīvi apvieno iepriekšējo slāņu pazīmes, lai iegūtu arvien
sarežģītākas pazīmes. Ieejas pazīmju karte tiek sadalīta mazākās neironu apakškopās. Katrs nākamā slāņa neirons saņem
impulsus no savas neironu apakškopas un tie tiek reizināti ar svaru kodolu \cite{tavanaei2017}:
\begin{equation}
  x^{t,l}_{i,j} = \sum^{K}_{m=0} \sum^{K}_{n=0} w^l_{m,n} s^{t,l-1}_{i+m,j+n}
  \label{eq:lif_conv}
\end{equation}
\begin{equation}
  {u}^{t,l}_{i,j} = u^{t-1,l}_{i,j}(1 - s^{t,l}_{i,j}) + x^{t,l}_{i,j}
  \label{eq:lif_v}
\end{equation}
\begin{equation}
  s^{t,l}_{i,j} = \Theta ({u}^{t,l}_{i,j} - \vartheta_{th}),
  \label{eq:lif_spike}
\end{equation}
kur \(w^l_{m,n}\) ir svaru kodols, kas ir kopīgs visam \(l\)-tajam slānim, lai samazinātu parametru skaitu līdzīgi kā tradicionālajos \Gls{cnn}, pieņemot, ka vienā pozīcijā noderīga kodola pazīme ir noderīga arī citās telpiskās pozīcijās. \(K^2\) ir ieejas pazīmju kartes kodola neironu skaits, ko izmanto \(i\)-tā neirona ieejas strāvas \(x^{t,l}_{i,j}\) ģenerēšanai, un \(s^{t,l-1}_{i+m,j+n}\) ir impulsi no \(i\)-tā ieejas kodola laika solī \(t+1\).

Tomēr praktiskos gadījumos ieejas tenzoriem ir vairāki ieejas kanāli \(C_{in}\) (piemēram, RGB krāsu kanāli vai pazīmju karšu kanāli). Tad katrs no \(C_{out}\) izejas kanāliem \(c\) tiek iegūts, veicot ieejas konvolūciju ar atsevišķu 3D kodolu, kura izmērs ir \(C_{in} \times K \times K\), un sinaptisko strāvu telpiskajā pozīcijā \(i,j\) katrā izejas kanālā nosaka
\begin{equation}
    x^{t,l}_{c,i,j} = \sum^{C_{in}}_{d=0} \sum^{K}_{m=0} \sum^{K}_{n=0} w^l_{c,d,m,n} s^{t,l-1}_{d,i+m,j+n},
    \label{eq:channel_conv}
\end{equation}
kur \(d\) indeksē ieejas kanālus, \(c\) indeksē izejas kanālus un \(w^l_{c,d,m,n}\) ir koplietotie konvolūcijas svari slānī \(l\). Līdzīgi katram izejas kanālam tiek ģenerēts membrānas potenciāls \({u}^{t,l}_{c,i,j}\) un impulsi \(s^{t,l}_{c,i,j}\).

Šādi \((K \times K)\) kodolu kopu izmanto \((H \times W)\) pikseļu ieejas pazīmju kartes konvolūcijai, iegūstot izejas pazīmju kartes ar
\begin{equation}
  \left( \frac{H - K + 2p}{s} + 1 \right) \times \left( \frac{W - K + 2p}{s} + 1 \right)
\end{equation}
neironiem, kur \(p\) ir papildinājuma apjoms pie malām un \(s\) ir solis - pikseļu skaits, par kuru kodolu pārvieto vienā reizē.

\begin{figure}[h!]
\begin{adjustbox}{max width=\linewidth,center}
\centering
    \begin{tikzpicture}[
        >=Stealth,
        neuron/.style={circle, draw, fill=gray!50, minimum size=0.9cm},
        spike/.style={rectangle, draw, minimum width=1.2cm, minimum height=0.6cm},
        weight/.style={midway, fill=white, inner sep=1pt},
        kcell/.style={draw, minimum width=0.6cm, minimum height=0.6cm, font=\scriptsize, fill=orange!15},
    ]
    \draw[step=0.6cm, gray] (0,0) grid (4.2,4.2);
    \foreach \x in {1.2,1.8,2.4}{
        \foreach \y in {1.2,1.8,2.4}{
            \filldraw[fill=gray!30, draw=black] (\x,\y) rectangle ++(0.6,0.6);
        }
    }
    \foreach \i in {0,...,6}{
        \foreach \j in {0,...,6}{
            \pgfmathtruncatemacro{\spike}{rnd>0.75}
            \node at ({0.3+0.6*\i},{0.3+0.6*\j}) {\spike};
        }
    }
    \node[below=0.2cm] at (2.1,0) {\small Impulsu pazīmju karte $s^{t,l-1}$};
    \node[above=0.15cm] at (2.1,4.2) {\small $3 \times 3$ uztveres lauks};

    \foreach \r in {0,...,2} {
        \foreach \c in {0,...,2} {
            \pgfmathtruncatemacro{\rr}{\r+1}
            \pgfmathtruncatemacro{\cc}{\c+1}
            \node[kcell] (k-\r-\c) at ({6.4+0.8*\c},{2.1+0.6*(1-\r)}) {$w_{\rr\cc}$};
        }
    }
    \node[above=0.15cm] at (6.7,2.1+0.6+0.3) {\small Kodols $w^l_{m,n}$};

    \node at (5.15,2.1) {\Large $\circledast$};
    \draw[-] (3.0,1.8) -- (4.65,2.1);
    \draw[-] (3.0,2.4) -- (4.65,2.1);

    \node[neuron, right=1.4cm of k-1-2] (sum) {$\sum_{m,n}$};
    \draw[->] (k-0-2) -- (sum);
    \draw[->] (k-2-2) -- (sum);

    \node[neuron, right=1.8cm of sum] (spikef) {${u}^{t,l}_{i,j}$};
    \node[spike, right=1.8cm of spikef] (y) {$s^{t,l}_{i,j}$};
    \draw[->] (sum) -- (spikef);
    \draw[->] (spikef) -- node[weight] {$\Theta(\cdot)$} (y);

    \node[above=0.3cm] at (9.6,2.6) {\small Konvolūcija};
    \node[above=0.3cm of spikef] {\small Impulsu neirons};
    \end{tikzpicture}
\end{adjustbox}
\caption{Impulsu konvolūcijas slānis. Koplietots reālu vērtību $3 \times 3$ kodols $w^l_{m,n}$ tiek konvolvēts ar binārās ieejas impulsu pazīmju kartes $s^{t,l-1}$ uztveres lauku. Svērto summu integrē impulsu neirons, kas ģenerē izejas impulsu $s^{t,l}_{i,j}$ pozīcijā $(i,j)$.}
\label{fig:spiking_conv_receptive_field}
\end{figure}


